\chapter{Implementación}
\label{chap:implementacion}

\lettrine{L}{a} implementación de PadelCenter se divide en dos proyectos
independientes que se comunican a través del contrato OpenAPI: el backend
en Java~21 con Spring Boot~3.4.4 y el frontend en TypeScript con Next.js~14.
En este capítulo se describen las decisiones de implementación más relevantes
de cada uno y el flujo de integración end-to-end.

\section{Backend}
\label{sec:implementacion_backend}

\subsection{Estructura del proyecto}

El proyecto backend organiza el código siguiendo la estructura de la arquitectura
hexagonal, con el paquete raíz \texttt{com.bookings.padelcenter}:

\begin{itemize}
  \item \texttt{domain.model}: records inmutables de dominio (User, Center, Field,
    Booking, Tournament, TournamentPair, Match, CenterRole, Auditable\ldots).
    Sin dependencias externas.
  \item \texttt{domain.port}: interfaces de puertos de entrada y salida.
  \item \texttt{domain.exception}: excepciones de dominio tipadas
    (ResourceNotFoundException, BookingAlreadyCancelledException,
    EmailAlreadyExistsException, UnauthorizedException\ldots).
  \item \texttt{application.create/read/update/delete}: 42 casos de uso
    organizados por naturaleza de la operación. Cada uno implementa
    \texttt{CommandUseCase} o \texttt{QueryUseCase}.
  \item \texttt{infrastructure.inbound.web}: controladores HTTP que implementan
    las interfaces generadas por OpenAPI.
  \item \texttt{infrastructure.inbound.mapper}: mappers API↔dominio.
  \item \texttt{infra\-structure.out\-bound.db.entity}: entidades JPA.
  \item \texttt{infra\-structure.out\-bound.db.mapper}: mappers dominio↔entidad JPA.
  \item \texttt{infra\-structure.out\-bound.security}: JWT, BCrypt, evaluadores
    de autorización.
  \item \texttt{infrastructure.config}: configuración de Spring Security y logging.
\end{itemize}

\subsection{Generación de código OpenAPI}

El plugin \texttt{openapi-generator-maven-plugin}~\cite{openapiGenerator} (v7.6.0) se configura en el
\texttt{pom.xml} para generar en fase \texttt{generate-sources}:

\begin{itemize}
  \item \textbf{Interfaces de controlador}: anotadas con \texttt{@RequestMapping},
    que los controladores implementan. Si la firma no coincide con el contrato,
    el compilador lo detecta.
  \item \textbf{DTOs como records Java}: inmutables, con anotaciones Bean
    Validation derivadas del esquema OpenAPI (\texttt{@NotNull}, \texttt{@Size},
    etc.). Exclusivos de la capa web; nunca se exponen al dominio.
\end{itemize}

\subsection{Patrón de mapeo en dos capas}

El proyecto implementa el mapeo entre capas mediante dos familias de clases
\texttt{@Component} manuales, sin dependencias de generación de código:

\begin{description}
  \item[\texttt{infrastructure.inbound.mapper}:] Mappers \textbf{API↔dominio}.
    Transforman los records DTO generados por OpenAPI en objetos de dominio
    (comandos y queries) antes de invocar el caso de uso, y transforman el
    resultado del caso de uso en el DTO de respuesta. Por ejemplo,
    \texttt{BookingApiMapper} convierte \texttt{CreateBookingRequest} en
    \texttt{CreateBookingUseCase.Command} y \texttt{Booking} en
    \texttt{BookingResponse}.

  \item[\texttt{infra\-structure.out\-bound.db.mapper}:] Mappers
    \textbf{dominio↔entidad JPA}. Transforman los records de dominio en
    entidades anotadas con \texttt{@Entity} y viceversa, gestionando la
    conversión de tipos (UUID, BigDecimal, enumerados) y la hidratación
    de objetos relacionados. Por ejemplo, \texttt{BookingPersistenceMapper}
    transforma entre \texttt{Booking} (record inmutable) y
    \texttt{BookingEntity} (entidad JPA mutable).
\end{description}

Esta separación garantiza que ni los DTOs de la API ni las entidades JPA
contaminan el dominio o los casos de uso.

\subsection{Flyway y migraciones de base de datos}

El esquema se gestiona con Flyway~\cite{flyway}. Las migraciones viven en
{\ttfamily src/\-main/\-resources/\-db/\-migration/}:

\begin{itemize}
  \item \texttt{V1\_\_init\_schema.sql}: tablas \texttt{users}, \texttt{centers},
    \texttt{fields}, \texttt{bookings}, \texttt{booking\_participants} y
    \texttt{center\_roles} con columnas de auditoría.
  \item \texttt{V3\_\_add\_missing\_indexes.sql}: índices de optimización.
  \item \texttt{V4\_\_add\_booking\_overlap\_index.sql}: índice parcial para
    detección eficiente de solapamiento de reservas a nivel de base de datos.
  \item \texttt{V5\_\_create\_tournament\_tables.sql}: crea las tablas de torneos
    (\texttt{tournaments}, \texttt{tournament\_pairs} y \texttt{tournament\_matches}).
  \item \texttt{V8\_\_add\_team\_name\_to\_pairs.sql}: campo
    \texttt{team\_name} opcional en \texttt{tournament\_pairs}.
\end{itemize}

Además de las migraciones de esquema, el perfil \texttt{local} carga una
ubicación adicional, \texttt{src/main/resources/db/seed/}, con datos de
prueba para desarrollo manual:

\begin{itemize}
  \item \texttt{V9\_\_seed\_curated\_data.sql}: conjunto de datos curado a
    mano (un usuario ADMIN, 6 centros, 24 pistas, un torneo con 8 parejas
    confirmadas, etc.) con identificadores UUID fijos, de forma que sigan
    siendo válidos tras un \texttt{docker compose down -v} que recree la
    base de datos desde cero.
  \item \texttt{V10\_\_seed\_bulk\_data.sql}: volumen adicional (180
    usuarios y \textasciitilde1000 reservas) generado con
    \texttt{generate\_series} de PostgreSQL, para probar paginación y
    filtros con datos realistas.
\end{itemize}

Esta ubicación solo se añade a \texttt{spring.flyway.locations} cuando el
perfil activo es \texttt{local} (ver \texttt{application-local.yaml}), por
lo que nunca se aplica en el perfil \texttt{test} ni en un despliegue real.

\subsection{Virtual threads (Project Loom)}

El servidor usa virtual threads de Java~21~\cite{projectloom}, activados con una única propiedad:

\begin{lstlisting}
spring.threads.virtual.enabled=true
\end{lstlisting}

Los virtual threads permiten un modelo de concurrencia de alta densidad: como
son baratos de crear y bloquear, el servidor puede manejar miles de peticiones
simultáneas sin el overhead de un pool de hilos de plataforma.

\subsection{Trazabilidad con MDC y logging estructurado}

Todas las peticiones HTTP pasan por un filtro \texttt{MdcFilter} que inyecta
en el MDC (\textit{Mapped Diagnostic Context}) el identificador de usuario y
un ID de correlación de petición. Los logs se emiten en formato JSON mediante
\textit{logstash-logback-encoder}, lo que facilita la ingesta en sistemas de
observabilidad como Elasticsearch o Loki.

\subsection{Convenciones de los casos de uso}

Cada caso de uso implementa \texttt{CommandUseCase<C, R>} o
\texttt{QueryUseCase<Q, R>}, donde el tipo genérico de entrada es siempre un
record interno de la propia clase. Esto agrupa comando y caso de uso en un
único fichero y hace explícito el contrato de cada operación. Las excepciones
de dominio se propagan hasta el \texttt{GlobalExceptionHandler}
(\texttt{@ControllerAdvice}), que las traduce a respuestas \texttt{ProblemDetail}
según RFC~9457.

\section{Frontend}
\label{sec:implementacion_frontend}

\subsection{Estructura del proyecto}

El frontend usa el App Router de Next.js~14~\cite{nextjs}. Las rutas se organizan en grupos
entre paréntesis para compartir layouts sin añadir segmentos a la URL:

\begin{itemize}
  \item \texttt{app/(auth)/}: rutas autenticadas con layout de aplicación
    (header, navegación). Incluye: \texttt{dashboard}, \texttt{centers},
    \texttt{centers/[centerId]}, \texttt{bookings}, \texttt{tournaments},
    \texttt{tournaments/[tournamentId]}, \texttt{matches} y \texttt{profile}.
  \item \texttt{app/(admin)/admin/}: rutas exclusivas de administrador.
    Incluye: gestión de \texttt{centers}, \texttt{centers/[centerId]},
    \texttt{fields}, \texttt{bookings} y \texttt{tournaments}.
  \item Rutas públicas: \texttt{/login}, \texttt{/register},
    \texttt{/auth/callback}.
\end{itemize}

\subsection{Librería de componentes: shadcn/ui y Radix UI}

La interfaz se construye sobre \textbf{shadcn/ui}~\cite{shadcnui}, una colección
de componentes copiables basada en primitivos accesibles de \textbf{Radix UI}~\cite{radixui} y estilizada con
Tailwind CSS~\cite{tailwind}. Esta elección aporta:

\begin{itemize}
  \item Componentes accesibles (ARIA) sin configuración adicional
    (Dialog, Popover, DropdownMenu, Select, Tabs, Toast\ldots).
  \item Control total sobre el código: los componentes se copian al repositorio
    y son modificables, a diferencia de librerías opacas de terceros.
  \item Integración nativa con Tailwind CSS y \texttt{class-variance-authority}
    para variantes tipadas de estilos.
\end{itemize}

\subsection{Generación de hooks con Orval}

Orval~\cite{orval} lee la especificación OpenAPI del backend corriendo en \texttt{:8080} y
genera en \texttt{src/lib/api/generated/}, organizado por tag (auth, bookings,
centers, fields, matches, tournaments, users):

\begin{itemize}
  \item Un hook React Query~\cite{reactquery} (\texttt{useQuery}) por operación GET.
  \item Una función de mutación (\texttt{useMutation}) por operación
    POST/PUT/DELETE.
  \item Los tipos TypeScript correspondientes a todos los esquemas del OpenAPI.
\end{itemize}

El cliente HTTP subyacente es Axios~\cite{axios} con un mutator personalizado que gestiona
el refresco automático del token cuando el servidor responde con HTTP~401.

\subsection{Gestión del estado de autenticación con Zustand}

\begin{sloppypar}
El \textit{access token} JWT se almacena en memoria en un store
Zustand~\cite{zustand} (\texttt{src/lib/auth/store.ts}), nunca en
\texttt{localStorage} ni en \texttt{sessionStorage}, lo que elimina el riesgo
de robo mediante XSS. Los datos del usuario se persisten en \texttt{localStorage}
mediante el middleware \texttt{persist} de Zustand para sobrevivir a recargas
de página. Al recargar, el interceptor de Axios intenta el refresco del token
mediante la cookie \textit{httpOnly} antes de redirigir al login.
\end{sloppypar}

\subsection{Protección de rutas}

El middleware de Next.js (\texttt{middleware.ts}) actúa como guardián de rutas:
verifica la autenticación para las rutas bajo \texttt{/(auth)/} y además
comprueba el rol \texttt{ADMIN} para las rutas bajo \texttt{/(admin)/},
redirigiendo a \texttt{/dashboard} si el rol es insuficiente.

\section{Flujo end-to-end: creación de una reserva}
\label{sec:flujo_end_to_end}

A continuación se describe el flujo completo desde el clic del usuario hasta
la respuesta en pantalla:

\begin{sloppypar}
\begin{enumerate}
  \item El usuario selecciona pista, fecha y franja horaria y pulsa
    \textit{Reservar}.
  \item El hook \texttt{useCreateBooking} (generado por Orval) envía
    \texttt{POST /api/v1/bookings} con el \textit{access token} JWT en el
    header \texttt{Authorization: Bearer}.
  \item El filtro de Spring Security decodifica el token con
    \texttt{NimbusJwtDecoder} y establece el contexto de seguridad.
  \item \texttt{BookingController} recibe la petición, la valida con Bean
    Validation, usa \texttt{BookingApiMapper} para transformar el DTO en un
    \texttt{CreateBookingUseCase.Command} y ejecuta el caso de uso.
  \item \texttt{CreateBookingUseCase} invoca el método
    \texttt{existsConflict()} de \texttt{BookingRepository}, que consulta el
    índice de solapamiento de la base de datos.
  \item Si no hay conflicto, persiste la reserva con estado \texttt{PENDING},
    calcula el precio total y devuelve el record \texttt{Booking}.
  \item El controlador usa \texttt{BookingApiMapper} para transformar el record
    en \texttt{BookingResponse} y responde con HTTP~201.
  \item React Query invalida la caché de \texttt{useGetMyBookings} y el listado
    se actualiza automáticamente.
\end{enumerate}
\end{sloppypar}

\section{Entorno local con Docker Compose}
\label{sec:entorno_local}

El fichero \texttt{tools/docker/docker-compose.yaml} define los servicios
necesarios para el desarrollo local:

\begin{itemize}
  \item \textbf{postgres}: PostgreSQL~16, expuesto en el puerto 5432. Flyway
    aplica las migraciones automáticamente al arrancar el backend.
  \item \textbf{padelcenter}: backend Spring Boot con perfil de desarrollo.
\end{itemize}

El frontend se ejecuta fuera de Docker con \texttt{pnpm dev} para aprovechar
el \textit{hot reload} de Next.js. La variable de entorno
\texttt{NEXT\_PUBLIC\_API\_URL=http://localhost:8080} conecta ambos proyectos.

\subsection{Colección Postman}

\texttt{tools/postman/} incluye una colección con las 48 operaciones de la
API, organizada por recurso, y un \textit{environment} (\texttt{PadelCenter
— Local}) con las credenciales y los identificadores fijos de los datos
curados de \texttt{V9}. Cada petición que crea un recurso captura su
identificador en una variable del entorno mediante un script de test, de
forma que la colección puede recorrerse de principio a fin sin intervención
manual. La validación sistemática de las 48 peticiones con
\texttt{newman}, la CLI de Postman, contra el entorno local fue precisamente
el método que permitió detectar varios de los defectos corregidos durante
el desarrollo (véase la sección~\ref{sec:resumen_cobertura}).
